% ============================================
% Johns Hopkins Letter Template
% ============================================

\documentclass[11pt,letterpaper]{letter}

\usepackage{graphicx}
\usepackage{eso-pic}
\usepackage{geometry}
\usepackage{microtype}
\usepackage[utf8]{inputenc}
\usepackage[hidelinks]{hyperref}

\geometry{left=25mm,right=25mm,top=25mm,bottom=25mm}

\newcommand{\PutJHULogoFG}[1]{%
  \AddToShipoutPictureFG*{%
    \AtPageUpperLeft{%
      \hspace{10mm}%
      \raisebox{-38mm}{%
        \includegraphics[width=6cm]{#1}%
      }%
    }%
  }%
}

\signature{Decory Edwards}
\address{
Department of Economics \\
Johns Hopkins University \\
Baltimore, MD 21218 \\
\texttt{dedwar65@jh.edu}
}

\begin{document}
\PutJHULogoFG{Images/jhulogo.png}

\begin{letter}{
Search Committee \\
Department of Economics \\
Hobart and William Smith Colleges \\
Geneva, NY
}

\opening{Dear Members of the Search Committee,}

I am writing to apply for the Visiting Assistant Professor position in Economics at Hobart and William Smith Colleges. I am a Ph.D.\ candidate in the Department of Economics at Johns Hopkins University (advisors: Christopher D.\ Carroll, Nicholas W.\ Papageorge, and Stelios Fourakis), and I expect to receive my degree in May 2026. My primary specializations are macroeconomic modeling, wealth and income dynamics, and computational economics. I am committed to high quality undergraduate teaching and inclusive pedagogy in liberal arts settings.

My research agenda focuses on the ability of macroeconomic models to generate empirically realistic distributions of wealth and income over the life cycle. A key driver of that work is modeling heterogeneity in the rate of return on assets, and understanding how return dispersion contributes to portfolio accumulation across households.

In my job market paper, I calibrate a life cycle heterogeneous agent model to match wealth moments from the Survey of Consumer Finances by allowing households to differ in their portfolio returns. The model helps explain observed wealth concentration and return dispersion. In a related research project using the Health and Retirement Study, I examine how trust influences both income growth and asset returns. I find a hump shaped relationship for trust and returns (consistent with the literature) and characterize the distribution of the persistent component of returns to net worth. These experiences have equipped me to frame research strategies and design modeling approaches, execute econometric and simulation analysis with large micro data sets, and translate results into clear and rigorous written deliverables for both academic and practitioner audiences.

I am especially enthusiastic about teaching intermediate macroeconomics and financial macro topics in a liberal arts environment. My research in heterogeneous agent macroeconomics provides a strong foundation for discussing business cycles, asset markets, and monetary and fiscal policy in ways that connect theory to contemporary financial developments. I am also prepared and willing to teach undergraduate statistics or econometrics, and I view quantitative reasoning as central to helping students evaluate economic evidence critically. I would welcome the opportunity to contribute to the general education curriculum and to design electives that connect macroeconomic theory to questions about financial markets, inequality, and policy.

Hobart and William Smith’s emphasis on mentorship, intellectual engagement, and preparing students for meaningful lives resonates strongly with me. I strive to create classrooms that are rigorous yet supportive, where students from diverse backgrounds feel encouraged to participate and develop confidence in their analytical skills. I am enthusiastic about relocating to Geneva and engaging fully in the residential academic community of the Finger Lakes region. In a five or six course teaching load, I would be eager to contribute actively to both the department and the broader campus community.

I believe I would be a strong fit for this role because:
\begin{enumerate}
    \item My specialization in macroeconomics and financial economics aligns directly with the department’s stated needs.
    \item I am committed to inclusive and discussion driven undergraduate teaching.
    \item I bring strong quantitative training that supports instruction in statistics and econometrics.
\end{enumerate}

Thank you very much for your consideration. I would welcome the opportunity to contribute to the Department of Economics at Hobart and William Smith Colleges.

\closing{Sincerely,}

\end{letter}
\end{document}
